%!TEX root = newmain.tex

{\Stipula} is a domain-specific programming language designed to model
legal contracts with formally enforceable properties, particularly
those governing asset transfers, obligations, and time-dependent
rights~\cite{CrafaLSV23}.  Contracts written in {\Stipula} combine
declarative clauses with executable semantics, yielding models whose
behaviour evolves dynamically over time. To facilitate the contract
lifecycle, a supporting toolchain has been developed that has been
integrated into a unified workbench, enabling legal practitioners and
developers in writing, testing, debugging, and managing legal
contracts~\cite{Laneve2026,stipulaprototype}.  While this
infrastructure provides practical support for contract development and
experimentation, it currently lacks advanced deductive verification
capabilities.

This paper complements the existing {\Stipula} workbench with a formal
verification layer capable of establishing both partial and total
correctness properties of legal contracts.  Our approach relies on a
systematic translation of {\Stipula} contracts into {\Java} programs
annotated with {\tt Java Modeling Language} ({\JML}) specifications,
thereby reducing the verification of legal contracts to the discharge
of proof obligations within a mature and well-established formal
methods framework.  Concretely, we use {\KeY}~\cite{KeYTutorial24} as
a general-purpose deductive verification system.

A central challenge in this setting is achieving sufficient precision
in the representation of execution behaviour. Legal contracts
typically exhibit \emph{dynamic} features -- notably the scheduling of
events at future times -- while {\JML} specifications are inherently
\emph{static}.  Modelling time and deferred events within a deductive
verification framework, therefore, requires an explicit and carefully
structured representation of the contract's temporal evolution. Unlike
our preliminary study~\cite{iFM25}, where time and events were treated
in an abstract manner, the present work fully incorporates timed
events into the translation. This is challenging because the
underlying semantic models of {\Stipula} contracts are infinite-state
due to the unbounded progression of time and the generation of future
events.

To represent the behaviour of {\Stipula} using static method contracts, we exploit the 
finite-state structure defined by the contract's underlying automaton, in line 
with {\Stipula}'s state-based programming style. Starting from this automaton, 
we systematically construct all sequences of clauses that are legally admissible 
according to the contract specification. This construction yields an explicit 
representation of the contract's behavioural space, abstracting away from concrete 
executions while preserving the permitted control-flow structures. A key component 
of this process is the dispatch table, which records the events that can be scheduled 
together with their trigger times. This representation plays a central role in 
verification: the soundness and completeness of the properties that can be established depend 
directly on how precisely the admissible execution paths are captured.

We first present the technique for acyclic {\Stipula} contracts. In
this setting, we formally define the translation, describe the
construction of the dispatch table and the generation of execution
paths, and show how partial-correctness properties of contracts can be
reduced to proof obligations discharged automatically by {\KeY}. For
acyclic contracts, the behavioural space is finite and can be
represented with high precision.

We subsequently extend our approach to cover cyclic contracts. The
translation devised for the acyclic setting can be generalised to
contracts featuring \emph{separate cycles}: cycles are not
nested, even if multiple paths between their entry/exit states
are permitted. Here, each cycle is mapped to a bounded {\Java} loop
iteration, whose termination is enforced through an explicit counter
and a predetermined upper bound.  This encoding makes the verification
task possible for a deductive verification tool, while
deliberately introducing a controlled over-approximation of the
contract's potentially unbounded behaviour. As customary in deductive
verification, the presence of loop iterations requires suitable
invariants. Accordingly, while the structural aspects of the
translation are fully automated, the user may be required to provide
contract-specific invariants to enable {\KeY} to complete the proof
obligations generated by cyclic constructs.

Overall, the main contributions of this work are as follows:
%
\begin{itemize}
\item \emph{A principled encoding of time-dependent contractual
    behaviour.}  We devise a novel and adequate data structure, the
  \emph{dispatch table}, for representing time constraints and event
  scheduling within a deductive verification setting.
\item \emph{A systematic and extensible translation framework.}  We
  define a formal, fully specified translation from legal contracts
  written in {\Stipula} into {\Java} programs annotated with {\JML}
  specifications. The translation is not limited to acyclic
  interaction patterns, but also covers separate (non-nested) cycles, 
  thereby addressing a wide range of
  realistic contractual behaviour.
\item \emph{A sound verification methodology based on mainstream
    deductive technology.}  We establish the soundness of the
  translation, ensuring that properties verified on the generated
  {\JML}/{\Java} code by a state-of-the-art deductive verification
  system such as {\KeY} correctly reflect the semantics of the
  original contract.
\end{itemize}

The paper is organised as follows. Section~\ref{sec:background}
contains the necessary background about {\Stipula} and
{\KeY}. Section~\ref{sec:translation} defines the translation function
for {\Stipula} contracts; the translation covers both acyclic
contracts and those with cycles.  Section~\ref{sec:dispatchtable}
analyses the management of timed events with the dispatch table and
the generation of execution paths; Section~\ref{sec:cyclic-behavior}
extends the analysis to cyclic behaviour.
Section~\ref{sec:relatedworks} briefly reviews related efforts in
formal verification of legal contracts and
Section~\ref{sec:conclusion} concludes the paper and outlines
directions for future work.




%%% Local Variables:
%%% mode: latex
%%% TeX-master: "newmain"
%%% End:
